\documentclass[11pt]{article}
\usepackage{amsmath,amssymb}

\title{Contrastive Self-Supervised Pretraining without Negative Pairs}
\author{C. Aydin \and B. Hollander \and S. Park}
\date{}

\begin{document}
\maketitle

\begin{abstract}
Most contrastive frameworks for visual representation learning rely on a
large pool of negative samples, requiring either huge batches or a memory
bank of stale features. We show that a stop-gradient asymmetric Siamese
network, trained only on positive pairs, matches or exceeds the linear-probe
accuracy of much larger contrastive baselines on ImageNet. The key
ingredient is the asymmetry itself: a predictor MLP attached to the online
network and a stop-gradient on the target.
\end{abstract}

\section{Setup}
Given an image $x$, two augmented views $x_1, x_2$ are produced by a stochastic
augmentation $\mathcal{T}$. An encoder $f$ maps each view to a representation,
and a small projector $g$ embeds it. The asymmetric branch additionally applies
a predictor $h$:
\begin{align}
p_1 &= h(g(f(x_1))), \\
z_2 &= \mathrm{sg}(g(f(x_2))),
\end{align}
where $\mathrm{sg}(\cdot)$ denotes the stop-gradient operator.

\section{Objective}
We minimise the negative cosine similarity, symmetrised over the two views:
\begin{equation}
\mathcal{L} = -\frac{1}{2}\!\left( \frac{p_1 \cdot z_2}{\|p_1\| \|z_2\|}
                                  + \frac{p_2 \cdot z_1}{\|p_2\| \|z_1\|}
                            \right).
\end{equation}
There are no explicit negatives, no temperature, no memory bank. Collapse to
a constant representation is prevented entirely by the asymmetry: gradients
flow through the predictor branch but not the target branch, and the
predictor cannot reach the trivial solution because the target it must match
is itself a non-stationary moving function of the input.

\section{Why doesn't it collapse?}
We analyse the gradient near the collapsed solution. Linearising around
$f \equiv c$ for constant $c$, the predictor $h$ must learn a non-trivial
mapping to compensate for the target's variance under augmentation. As long
as $h$ has finite capacity and is updated by gradient descent, the system
has a non-trivial fixed point. In practice, removing the predictor or the
stop-gradient causes immediate collapse within a few epochs.

\section{Augmentations}
We use a strong augmentation policy: random resized crop, colour jitter,
Gaussian blur, horizontal flip. Removing colour jitter is the most damaging
single ablation, costing $9\%$ linear-probe accuracy. The intuition is that
without colour jitter, the network can solve the task by learning a
colour-histogram fingerprint that ignores higher-level content.

\section{Results}
\begin{tabular}{lcc}
\hline
Method & Batch & Linear probe \\
\hline
SimCLR     & 4096 & 69.3\% \\
MoCo v2    & 256  & 71.1\% \\
BYOL       & 4096 & 74.3\% \\
SimSiam (ours) & 256 & 71.3\% \\
\hline
\end{tabular}

The headline is not raw accuracy but compute: matching MoCo at the same batch
size with substantially fewer moving parts (no momentum encoder, no negative
queue) suggests the negative-pool machinery may be doing less work than the
literature implied.

\section{Transfer}
We fine-tune on PASCAL VOC detection and COCO instance segmentation. Pretraining
gains transfer cleanly: $+1.4$ AP on VOC, $+0.9$ on COCO over a supervised
ImageNet initialisation. Self-supervised features are now the better starting
point for most downstream tasks where label noise dominates.

\section{Discussion}
The contrastive learning literature has been slowly converging on the view
that "negative pairs" are a sufficient but not necessary anti-collapse
mechanism. The deeper question --- what minimum structural assumption forces
representations to be informative --- remains open. Our results put another
data point on the side of architectural rather than loss-based regularisation.

\bibliographystyle{plain}
\begin{thebibliography}{9}
\bibitem{chen2021exploring} X. Chen, K. He. \emph{Exploring Simple Siamese Representation Learning.} CVPR, 2021.
\bibitem{grill2020bootstrap} J.-B. Grill et al. \emph{Bootstrap Your Own Latent.} NeurIPS, 2020.
\end{thebibliography}

\end{document}
